\par\Needspace{6\baselineskip}\medskip
\hypertarget{native-102-article077}{}\pdfbookmark[2]{77.}{article-077}
\providecommand{\DailyMotionSign}{\mathord{\text{\itshape 7}}}
\begin{center}77.\end{center}
Verbindet man die Herleitungen des vorigen Artikels mit dem, was wir in Art. 15, 16, 20, 27 und 28 für die einzelnen Arten von Kegelschnitten angegeben haben, so erhält man alles, was zur Berechnung der differenziellen Änderungen des geozentrischen Ortes infolge von Änderungen der einzelnen Elemente erforderlich ist. Zur weiteren Veranschaulichung dieser Regeln greifen wir das oben in Art. 13, 14, 51, 63 und 65 behandelte Beispiel wieder auf. Zunächst drücken wir nach dem vorigen Artikel $\mathrm{d}l$ und $\mathrm{d}b$ durch $\mathrm{d}r$, $\mathrm{d}u$, $\mathrm{d}i$ und $\mathrm{d}\AscendingNode$ aus. Die Rechnung verläuft wie folgt.
\begin{center}
\begin{tabular}{@{}ll@{}}
\makebox[17em][l]{$\log\operatorname{tang}\omega$\hfill .\,.\,.\,.} & $8{,}40113$\\[.4em]
\makebox[17em][l]{$\log\cos i$\hfill .\,.\,.\,.\,.} & $9{,}98853$\\[.4em]
\hline\noalign{\vskip .3em}
\makebox[17em][l]{$\log\operatorname{tang}M$\hfill .\,.\,.} & $8{,}41260$\\[.4em]
$M=$ & $1^0\,28'\,52''$\\[.4em]
$M-u=$ & $165\quad17\quad8$\\[.4em]
\end{tabular}
\end{center}
\begin{center}
\begin{tabular}{@{}ll@{}}
\makebox[17em][l]{$\log\sin\omega$\hfill .\,.\,.\,.\,.} & $8{,}40099_n$\\[.4em]
\makebox[17em][l]{$\log\operatorname{tang}i$\hfill .\,.\,.\,.} & $9{,}36723$\\[.4em]
\hline\noalign{\vskip .3em}
\makebox[17em][l]{$\log\operatorname{tang}N$\hfill .\,.\,.} & $7{,}76822_n$\\[.4em]
$N=$ & $179^0\,39'\,50''$\\[.4em]
$N-b=$ & $186\quad1\quad45$\\[.4em]
\end{tabular}
\end{center}
\begin{center}
\begin{tabular}{@{}ll@{}}
\makebox[17em][l]{$\log\operatorname{tang}(M-u)$\hfill .\,.\,.} & $9{,}41932_n$\\[.4em]
\makebox[17em][l]{$\log\cos\omega\sin i$\hfill .\,.\,.\,.\,.} & $9{,}35562_n$\\[.4em]
\hline\noalign{\vskip .3em}
\makebox[17em][l]{$\log\operatorname{tang}P$\hfill .\,.\,.\,.\,.\,.\,.\,.} & $0{,}06370$\\[.4em]
$P=$ & $49^0\,11'\,13''$\\[.4em]
$N-b-P=$ & $136\quad50\quad32$\\[.4em]
\end{tabular}
\end{center}
\SourceLeaf{103}
\begin{center}
\begin{tabular}{@{}ll@{}}
\multicolumn{2}{c}{I$^{*}$}\\[.4em]
\makebox[17em][l]{$\log\sin(L-l)$\hfill .\,.\,.} & $9{,}72125$\\[.4em]
\makebox[17em][l]{$\log R$\hfill .\,.\,.\,.\,.\,.\,.\,.\,.\,.} & $9{,}99810$\\[.4em]
\makebox[17em][l]{$\mathrm{C.}\log\Delta'$\hfill .\,.\,.\,.\,.\,.\,.} & $9{,}92027$\\[.4em]
\hline\noalign{\vskip .3em}
\makebox[17em][l]{$({}^{*})$\hfill .\,.\,.\,.\,.\,.\,.\,.\,.\,.\,.\,.} & $9{,}63962$\\[.4em]
\makebox[17em][l]{$\mathrm{C.}\log r$\hfill .\,.\,.\,.\,.\,.\,.\,.} & $9{,}67401$\\[.4em]
\hline\noalign{\vskip .3em}
\makebox[17em][l]{$\log\left(\frac{\mathrm{d}l}{\mathrm{d}r}\right)$\hfill .\,.\,.\,.\,.\,.\,.} & $9{,}31363$\\[.4em]
\end{tabular}
\end{center}
\begin{center}
\begin{tabular}{@{}ll@{}}
\multicolumn{2}{c}{II$^{**}$}\\[.4em]
\makebox[17em][l]{$({}^{*})$\hfill .\,.\,.\,.\,.\,.\,.\,.\,.\,.\,.} & $9{,}63962$\\[.4em]
\makebox[17em][l]{$\log\operatorname{cot}(M-u)$\hfill .} & $0{,}58068_n$\\[.4em]
\hline\noalign{\vskip .3em}
\makebox[17em][l]{$\log\left(\frac{\mathrm{d}l}{\mathrm{d}u}\right)$\hfill .\,.\,.\,.\,.\,.} & $0{,}22030$\\[.4em]
\end{tabular}
\end{center}
\begin{center}
\begin{tabular}{@{}ll@{}}
\multicolumn{2}{c}{III$^{*}$}\\[.4em]
\makebox[17em][l]{$\log\cos\omega$\hfill .\,.\,.\,.\,.\,.\,.\,.\,.\,.\,.} & $9{,}99986_n$\\[.4em]
\makebox[17em][l]{$\log\operatorname{tang}b$\hfill .\,.\,.\,.\,.\,.\,.\,.\,.\,.} & $9{,}04749_n$\\[.4em]
\hline\noalign{\vskip .3em}
\makebox[17em][l]{$\log\left(\frac{\mathrm{d}l}{\mathrm{d}i}\right)$\hfill .\,.\,.\,.\,.\,.\,.\,.\,.\,.\,.} & $9{,}04735_n$\\[.4em]
\end{tabular}
\end{center}
\begin{center}
\begin{tabular}{@{}ll@{}}
\multicolumn{2}{c}{IV}\\[.4em]
\makebox[17em][l]{$\log\frac{R}{\Delta'}$\hfill .\,.\,.\,.\,.\,.\,.\,.\,.\,.} & $9{,}91837$\\[.4em]
\makebox[17em][l]{$\log\cos(L-l)$\hfill .\,.} & $9{,}92956$\\[.4em]
\hline\noalign{\vskip .3em}
\makebox[17em][l]{$({}^{**})$\hfill .\,.\,.\,.\,.\,.\,.\,.\,.\,.\,.} & $9{,}84793$\\[.4em]
\multicolumn{2}{@{}l@{}}{$=\log\left(\frac{\mathrm{d}l}{\mathrm{d}\AscendingNode}-1\right)$}\\[.4em]
\end{tabular}
\end{center}
\begin{center}
\begin{tabular}{@{}ll@{}}
\multicolumn{2}{c}{V$^{*}$}\\[.4em]
\makebox[17em][l]{$({}^{**})$\hfill .\,.\,.\,.\,.\,.\,.\,.\,.\,.} & $9{,}84793$\\[.4em]
\makebox[17em][l]{$\log\sin b\cos b$\hfill .\,.} & $9{,}04212_n$\\[.4em]
\makebox[17em][l]{$\mathrm{C.}\log r$\hfill .\,.\,.\,.\,.\,.\,.} & $9{,}67401$\\[.4em]
\hline\noalign{\vskip .3em}
\makebox[17em][l]{$\log\left(\frac{\mathrm{d}b}{\mathrm{d}r}\right)$\hfill .\,.\,.\,.\,.\,.} & $8{,}56406$\\[.4em]
\end{tabular}
\end{center}
\begin{center}
\begin{tabular}{@{}ll@{}}
\multicolumn{2}{c}{VI$^{**}$}\\[.4em]
\makebox[17em][l]{$\log\frac{r}{\Delta}$\hfill .\,.\,.\,.\,.\,.\,.\,.\,.\,.\,.\,.\,.} & $0{,}24357$\\[.4em]
\makebox[17em][l]{$\log\sin(M-u)$\hfill .\,.\,.\,.\,.} & $9{,}40484$\\[.4em]
\makebox[17em][l]{$\log\cos(N-b-P)$\hfill .\,.} & $9{,}86301_n$\\[.4em]
\makebox[17em][l]{$\mathrm{C.}\log\sin P$\hfill .\,.\,.\,.\,.\,.\,.\,.} & $0{,}12099$\\[.4em]
\hline\noalign{\vskip .3em}
\makebox[17em][l]{$\log\left(\frac{\mathrm{d}b}{\mathrm{d}u}\right)$\hfill .\,.\,.\,.\,.\,.\,.\,.\,.\,.\,.} & $9{,}63241_n$\\[.4em]
\end{tabular}
\end{center}
\begin{center}
\begin{tabular}{@{}ll@{}}
\multicolumn{2}{c}{VII$^{*}$}\\[.4em]
\makebox[17em][l]{$\log r\sin u\cos i$\hfill .\,.} & $9{,}75999_n$\\[.4em]
\makebox[17em][l]{$\log\cos(N-b)$\hfill .\,.} & $9{,}99759_n$\\[.4em]
\makebox[17em][l]{$\mathrm{C.}\log\Delta$\hfill .\,.\,.\,.\,.\,.\,.} & $9{,}91759$\\[.4em]
\makebox[17em][l]{$\mathrm{C.}\log\cos N$\hfill .\,.\,.\,.} & $0{,}00001_n$\\[.4em]
\hline\noalign{\vskip .3em}
\makebox[17em][l]{$\log\left(\frac{\mathrm{d}b}{\mathrm{d}i}\right)$\hfill .\,.\,.\,.\,.\,.\,.} & $9{,}67518_n$\\[.4em]
\end{tabular}
\end{center}
\begin{center}
\begin{tabular}{@{}ll@{}}
\multicolumn{2}{c}{VIII}\\[.4em]
\makebox[17em][l]{$({}^{*})$\hfill .\,.\,.\,.\,.\,.\,.\,.\,.\,.\,.} & $9{,}63962$\\[.4em]
\makebox[17em][l]{$\log\sin b\cos b$\hfill .\,.} & $9{,}04212_n$\\[.4em]
\hline\noalign{\vskip .3em}
\makebox[17em][l]{$\log\left(\frac{\mathrm{d}b}{\mathrm{d}\AscendingNode}\right)$\hfill .\,.\,.\,.\,.\,.} & $8{,}68174_n$\\[.4em]
\end{tabular}
\end{center}
\par
Fasst man diese Werte zusammen, so ergibt sich
\[
\begin{aligned}
\mathrm{d}l={}&+0{,}20589\,\mathrm{d}r+1{,}66073\,\mathrm{d}u-0{,}11152\,\mathrm{d}i+1{,}70458\,\mathrm{d}\AscendingNode\\
\mathrm{d}b={}&+0{,}03665\,\mathrm{d}r-0{,}42895\,\mathrm{d}u-0{,}47335\,\mathrm{d}i-0{,}04806\,\mathrm{d}\AscendingNode.
\end{aligned}
\]
Wir brauchen hier wohl kaum zu wiederholen, worauf wir schon mehrfach hingewiesen haben. Entweder sind die Änderungen $\mathrm{d}l$, $\mathrm{d}b$, $\mathrm{d}u$, $\mathrm{d}i$ und $\mathrm{d}\AscendingNode$ in Teilen des Radius auszudrücken, oder die Koeffizienten von $\mathrm{d}r$ selbst sind mit $206265''$ zu multiplizieren, wenn man jene Änderungen in Bogensekunden angibt.
\par
Bezeichnen wir die Länge des Perihels (hier $52^0\,18'\,9''{,}30$) mit $\Pi$ und die wahre Anomalie mit $v$, so ist die Länge in der Bahn $=u+\AscendingNode=v+\Pi$. Also gilt $\mathrm{d}u=\mathrm{d}v+\mathrm{d}\Pi-\mathrm{d}\AscendingNode$. Setzt man diesen Wert in die vorigen Formeln ein, so erhält man $\mathrm{d}l$ und $\mathrm{d}b$ ausgedrückt durch $\mathrm{d}r$, $\mathrm{d}v$, $\mathrm{d}\Pi$, $\mathrm{d}\AscendingNode$ und $\mathrm{d}i$. \SourceLeaf{104}Es bleibt somit nur noch, $\mathrm{d}r$ und $\mathrm{d}v$ nach Art. 15 und 16 durch die differenziellen Änderungen der elliptischen Elemente auszudrücken\footnote{Jeder wird ohne Weiteres erkennen, dass das Zeichen $M$ in der folgenden Rechnung aufhört, unseren Hilfswinkel zu bezeichnen, und wieder wie in Abschnitt I für die mittlere Anomalie steht.}.
\par
In unserem Beispiel hatten wir in Art. 14 $\log\frac{r}{a}=9{,}90355=\log\left(\frac{\mathrm{d}r}{\mathrm{d}a}\right)$.
\begin{center}
\begin{tabular}{@{}ll@{}}
\makebox[17em][l]{$\log\frac{aa}{rr}$\hfill .\,.\,.\,.\,.\,.} & $0{,}19290$\\[.4em]
\makebox[17em][l]{$\log\cos\varphi$\hfill .\,.\,.\,.} & $9{,}98652$\\[.4em]
\hline\noalign{\vskip .3em}
\makebox[17em][l]{$\log\left(\frac{\mathrm{d}v}{\mathrm{d}M}\right)$\hfill .\,.\,.\,.} & $0{,}17942$\\[.4em]
\end{tabular}
\end{center}
\begin{center}
\begin{tabular}{@{}ll@{}}
$2-e\cos E=$ & $1{,}80085$\\[.4em]
$ee=$ & $0{,}06018$\\[.4em]
\hline\noalign{\vskip .3em}
 & $1{,}74067$\\[.4em]
\makebox[17em][l]{$\log$\hfill .\,.\,.\,.\,.\,.\,.\,.} & $0{,}24072$\\[.4em]
\makebox[17em][l]{$\log\frac{aa}{rr}$\hfill .\,.\,.\,.\,.\,.} & $0{,}19290$\\[.4em]
\makebox[17em][l]{$\log\sin E$\hfill .\,.\,.\,.} & $9{,}76634_n$\\[.4em]
\hline\noalign{\vskip .3em}
\makebox[17em][l]{$\log\left(\frac{\mathrm{d}v}{\mathrm{d}\varphi}\right)$\hfill .\,.\,.\,.\,.} & $0{,}19996_n$\\[.4em]
\end{tabular}
\end{center}
\begin{center}
\begin{tabular}{@{}ll@{}}
\makebox[17em][l]{$\log a$\hfill .\,.\,.\,.\,.\,.\,.} & $0{,}42244$\\[.4em]
\makebox[17em][l]{$\log\operatorname{tang}\varphi$\hfill .\,.\,.} & $9{,}40320$\\[.4em]
\makebox[17em][l]{$\log\sin v$\hfill .\,.\,.\,.\,.} & $9{,}84931_n$\\[.4em]
\hline\noalign{\vskip .3em}
\makebox[17em][l]{$\log\left(\frac{\mathrm{d}r}{\mathrm{d}M}\right)$\hfill .\,.\,.\,.} & $9{,}67495_n$\\[.4em]
\end{tabular}
\end{center}
\begin{center}
\begin{tabular}{@{}ll@{}}
\makebox[17em][l]{$\log a$\hfill .\,.\,.\,.\,.\,.\,.} & $0{,}42244$\\[.4em]
\makebox[17em][l]{$\log\cos\varphi$\hfill .\,.\,.\,.} & $9{,}98652$\\[.4em]
\makebox[17em][l]{$\log\cos v$\hfill .\,.\,.\,.} & $9{,}84966$\\[.4em]
\hline\noalign{\vskip .3em}
\makebox[17em][l]{$\log\left(\frac{\mathrm{d}r}{\mathrm{d}\varphi}\right)$\hfill .\,.\,.\,.\,.} & $0{,}25862_n$\\[.4em]
\end{tabular}
\end{center}
\noindent Daraus ergibt sich
\[
\begin{aligned}
\mathrm{d}v&=+1{,}51154\,\mathrm{d}M-1{,}58475\,\mathrm{d}\varphi\\
\mathrm{d}r&=-0{,}47310\,\mathrm{d}M-1{,}81393\,\mathrm{d}\varphi+0{,}80085\,\mathrm{d}a,
\end{aligned}
\]
und durch Einsetzen dieser Werte in die vorigen Formeln erhält man
\[
\begin{aligned}
\mathrm{d}l={}&+2{,}41287\,\mathrm{d}M-3{,}00531\,\mathrm{d}\varphi+0{,}16488\,\mathrm{d}a\\
 &+1{,}66073\,\mathrm{d}\Pi-0{,}11152\,\mathrm{d}i+0{,}04385\,\mathrm{d}\AscendingNode\\
\mathrm{d}b={}&-0{,}66572\,\mathrm{d}M+0{,}61331\,\mathrm{d}\varphi+0{,}02935\,\mathrm{d}a\\
 &-0{,}42895\,\mathrm{d}\Pi-0{,}47335\,\mathrm{d}i+0{,}38089\,\mathrm{d}\AscendingNode.
\end{aligned}
\]
\par
Es sei angenommen, dass der Zeitpunkt, zu dem der berechnete Ort gehört, $n$ Tage von der Epoche entfernt ist. Bezeichnet man die mittlere Länge zur Epoche mit $N$ und die tägliche Bewegung mit $\DailyMotionSign$, so gilt $M=N+n\DailyMotionSign-\Pi$, also $\mathrm{d}M=\mathrm{d}N+n\,\mathrm{d}\DailyMotionSign-\mathrm{d}\Pi$. In unserem Beispiel gehört der berechnete Ort zum $17{,}41501$. Oktober 1804, bezogen auf den Pariser Meridian. Nimmt man daher den Anfang des Jahres 1805 als Epoche, so ist $n=-74{,}58499$. Die für diese Epoche festgelegte mittlere Länge \SourceLeaf{105}war $=41^0\,52'\,21''{,}68$ und die tägliche Bewegung $=824''{,}7990$. Setzt man nun in die eben gefundenen Formeln den Wert für $\mathrm{d}M$ ein, so erhält man die folgenden differenziellen Änderungen des geozentrischen Ortes, ausgedrückt allein durch die Änderungen der Elemente.
\[
\begin{aligned}
\mathrm{d}l={}&+2{,}41287\,\mathrm{d}N-179{,}96\,\mathrm{d}\DailyMotionSign-0{,}75214\,\mathrm{d}\Pi\\
 &-3{,}00531\,\mathrm{d}\varphi+0{,}16488\,\mathrm{d}a-0{,}11152\,\mathrm{d}i+0{,}04385\,\mathrm{d}\AscendingNode\\
\mathrm{d}b={}&-0{,}66572\,\mathrm{d}N+49{,}65\,\mathrm{d}\DailyMotionSign+0{,}23677\,\mathrm{d}\Pi\\
 &+0{,}61331\,\mathrm{d}\varphi+0{,}02935\,\mathrm{d}a-0{,}47335\,\mathrm{d}i+0{,}38089\,\mathrm{d}\AscendingNode.
\end{aligned}
\]
\par
Wird die Masse des Himmelskörpers entweder vernachlässigt oder wenigstens als bekannt angesehen, so hängen $\DailyMotionSign$ und $a$ voneinander ab. Daher lässt sich entweder $\mathrm{d}\DailyMotionSign$ oder $\mathrm{d}a$ aus unseren Formeln eliminieren. Nach Art. 6 gilt $\DailyMotionSign a^{\frac{3}{2}}=k\BareRoot{(1+\mu)}$, also $\frac{\mathrm{d}\DailyMotionSign}{\DailyMotionSign}=-\frac{3}{2}\frac{\mathrm{d}a}{a}$. Soll in dieser Formel $\mathrm{d}\DailyMotionSign$ in Teilen des Radius ausgedrückt werden, so muss auch $\DailyMotionSign$ in derselben Weise ausgedrückt werden. Für unser Beispiel ergibt sich somit
\begin{center}
\begin{tabular}{@{}ll@{}}
\makebox[17em][l]{$\log\DailyMotionSign$\hfill .\,.\,.\,.\,.} & $2{,}91635$\\[.4em]
\makebox[17em][l]{$\log1''$\hfill .\,.\,.\,.} & $4{,}68557$\\[.4em]
\makebox[17em][l]{$\log\frac{3}{2}$\hfill .\,.\,.\,.\,.} & $0{,}17609$\\[.4em]
\makebox[17em][l]{$\mathrm{C.}\log a$\hfill .\,.\,.} & $9{,}57756$\\[.4em]
\hline\noalign{\vskip .3em}
\makebox[17em][l]{$\log\frac{\mathrm{d}\DailyMotionSign}{\mathrm{d}a}$\hfill .\,.\,.\,.} & $7{,}35557_n$\\[.4em]
\end{tabular}
\end{center}
\noindent oder $\mathrm{d}\DailyMotionSign=-0{,}0022676\,\mathrm{d}a$ und $\mathrm{d}a=-440{,}99\,\mathrm{d}\DailyMotionSign$. Setzt man diesen Wert in unsere Formeln ein, so erhält man schließlich ihre endgültige Form.
\[
\begin{aligned}
\mathrm{d}l={}&+2{,}41287\,\mathrm{d}N-252{,}67\,\mathrm{d}\DailyMotionSign-0{,}75214\,\mathrm{d}\Pi\\
 &-3{,}00531\,\mathrm{d}\varphi-0{,}11152\,\mathrm{d}i+0{,}04385\,\mathrm{d}\AscendingNode\\
\mathrm{d}b={}&-0{,}66572\,\mathrm{d}N+36{,}71\,\mathrm{d}\DailyMotionSign+0{,}23677\,\mathrm{d}\Pi\\
 &+0{,}61331\,\mathrm{d}\varphi-0{,}47335\,\mathrm{d}i+0{,}38089\,\mathrm{d}\AscendingNode.
\end{aligned}
\]
Bei der Herleitung dieser Formeln haben wir vorausgesetzt, dass alle Änderungen $\mathrm{d}l$, $\mathrm{d}b$, $\mathrm{d}N$, $\mathrm{d}\DailyMotionSign$, $\mathrm{d}\Pi$, $\mathrm{d}\varphi$, $\mathrm{d}i$ und $\mathrm{d}\AscendingNode$ in Teilen des Radius ausgedrückt sind. Wegen der Homogenität aller Glieder gelten dieselben Formeln aber offenbar auch dann noch, wenn alle diese Änderungen in Bogensekunden ausgedrückt werden.
\begin{center}\rule{7em}{.4pt}\end{center}
